%% ICASSP manuscript (Draft 11) on the OFFICIAL ICASSP 2027 template.
%% Reader-first scientific rewrite with a companion-repository layer for expanded numerical results,
%% provenance and reproduction. The abstract is unchanged from Draft 10.
%%
%% ICASSP 2027 allows four pages of technical content and an optional fifth page containing only
%% references, funding acknowledgements and an ethics statement.
%% draft11-paper-repo
\documentclass{article}
\usepackage{spconf,amsmath,amssymb,graphicx,booktabs,array,multirow,url,cite}

\newcommand{\Rn}{R_{\mathrm{nat}}}
\newcommand{\Rs}{R_{\mathrm{short}}}
\newcommand{\Rm}{R_{\mathrm{mus}}}
\newcommand{\ci}[2]{[{#1},\,{#2}]}
\newcommand{\PFT}{P{+}FT}
\emergencystretch=2em
\renewcommand{\topfraction}{0.92}
\renewcommand{\bottomfraction}{0.85}
\renewcommand{\textfraction}{0.06}
\renewcommand{\floatpagefraction}{0.75}
\setcounter{topnumber}{3}
\setcounter{totalnumber}{4}
\setcounter{dbltopnumber}{2}
\renewcommand{\dbltopfraction}{0.92}
\renewcommand{\dblfloatpagefraction}{0.75}
\setlength{\textfloatsep}{5pt plus 2pt minus 2pt}
\setlength{\floatsep}{6pt plus 2pt minus 2pt}
\setlength{\intextsep}{8pt plus 2pt minus 2pt}
\setlength{\dbltextfloatsep}{5pt plus 2pt minus 2pt}
\setlength{\dblfloatsep}{6pt plus 2pt minus 2pt}
\setlength{\abovecaptionskip}{3pt}
\setlength{\belowcaptionskip}{0pt}

\begin{document}
\ninept

\title{Recovery Fine-Tuning Recovers Where It Was Trained\\
Duration- and Domain-Dependent Gains in Pruned Text-to-Audio Diffusion}

\twoauthors
  {Gabriel Bibb\'o}
  {Independent researcher \\
   \texttt{gabobibbo@gmail.com}}
  {Arshdeep Singh, \, Mark D. Plumbley}
  {King's College London, UK \\
   \{\texttt{arshdeep.singh}, \texttt{mark.plumbley}\}\texttt{@kcl.ac.uk}}

\maketitle

\begin{abstract}
Structured pruning reduces the cost of text to audio diffusion models, and recovery fine tuning is commonly used to restore performance after pruning. Its success, however, is usually judged from one score at one inference setting, leaving open whether the recovered performance persists when generation conditions change. We evaluate released AudioLDM-M checkpoints at $65\%$ and $83\%$ pruning across clip duration and prompt domain. At $83\%$ pruning, fine tuning closes $63\%$ of the gap to real audio at the $10.24$\,s fine tuning duration, but only $33\%$ at $3.84$\,s, and yields no alignment gain on held out hip hop captions. The duration dependence replicates, grows monotonically across four durations, and persists across scorers, event level metrics, and the published sampler. These results establish operating point coverage as a necessary dimension of recovery evaluation, because a single setting can give an incomplete account of compressed model quality.
\end{abstract}

\begin{keywords}
text-to-audio generation, diffusion models, structured pruning, recovery fine-tuning, evaluation
\end{keywords}

\section{Introduction}
Text-to-audio generation has improved rapidly, but diffusion models remain expensive because the denoising network is evaluated repeatedly for every generated clip. AudioLDM-M, for example, contains a $416$\,M-parameter U-Net~\cite{liu2023audioldm}. Structured pruning offers a direct way to reduce this cost. Whole channels are removed, which reduces parameters and computation together. The resulting network usually loses quality, so pruning is followed by a recovery fine-tuning stage intended to regain the lost performance. This prune-then-recover pattern is established in image diffusion~\cite{fang2023diffpruning,kim2024bksdm,fang2025tinyfusion} and has recently been applied to text-to-audio generation~\cite{singh2026pruning}.

Recovery is often discussed as though it were a single property of the compressed checkpoint. In practice it is measured under a particular evaluation recipe. A benchmark fixes a prompt distribution, a requested clip duration, a sampler and a metric. If the recovered model is evaluated only under those conditions, a strong score shows that fine-tuning helped there. It does not show that the same improvement survives when the generation conditions change.

Two scientific claims are easy to conflate here. One concerns the quality of the final compressed model under the benchmark conditions. The other concerns what the recovery stage itself has restored and how broadly that restoration holds. The first can be answered by the final score of \PFT{}. The second requires comparison with P at more than one condition. Our study concerns this second claim. This distinction is useful whenever recovery is part of the compression pipeline because a good final checkpoint and a broadly transferable recovery effect are not the same object.

This distinction matters increasingly in audio generation. Recent evaluation work has moved beyond a single perceptual or alignment score. TTA-Bench, for example, evaluates text-to-audio systems across accuracy, robustness, generalization and other dimensions~\cite{wang2026ttabench}. Fine-grained benchmarks likewise expose capabilities that aggregate scores can hide. The compression literature, however, still largely reports post-recovery performance through fixed benchmark recipes. The missing question is therefore not whether recovery fine-tuning can improve a pruned model. It is whether the \emph{amount recovered} is stable across the conditions under which the model may actually be used.

We call those conditions the inference \emph{operating point}. Duration is one operating-point axis because AudioLDM is trained at a fixed $10.24$\,s clip length~\cite{liu2023audioldm,liu2024audioldm2}. Prompt domain is another because the recovery stage is performed on AudioCaps, and fine-tuning can improve in-distribution performance while degrading behaviour under distribution shift~\cite{kumar2022finetuning}. Sampler configuration provides a third axis. This framing turns recovery from one benchmark number into a function of how the model is queried.

We study this question using the released pruned and recovery-fine-tuned AudioLDM-M-Full checkpoints of Singh et al.~\cite{singh2026pruning}. We train no new model. Instead, we compare each pruned checkpoint P with the corresponding recovered checkpoint \PFT{} at matched operating points. The paper makes three contributions:
\begin{itemize}\itemsep1.5pt \parskip0pt \topsep2pt
\item We show that recovery depends strongly on requested duration. At $83\%$ pruning, the paired gain grows monotonically from $3.84$ to $10.24$\,s and is largest at the duration used for recovery fine-tuning.
\item We show that recovery is also domain-dependent. The same checkpoint improves strongly on AudioCaps while giving no text-audio alignment gain on held-out hip-hop captions.
\item We provide an evaluation protocol that separates recovery from absolute model quality using paired noise, chance and real-audio anchors, recovery ratios and pre-specified operating-point contrasts.
\end{itemize}
The manuscript is self-contained. Expanded numerical tables, exact intervals, prompt manifests, analysis status and reproduction scripts are indexed by paper section in the companion repository~\cite{bibbo2026repo}.

\section{Background and motivation}
\subsection{Recovery after structured pruning}
Structured pruning removes complete channels or layers rather than individual weights, which can translate model compression into real inference savings. Diffusion pruning methods therefore combine a pruning rule with a recovery stage that adapts the reduced network after capacity has been removed~\cite{fang2023diffpruning,kim2024bksdm,fang2025tinyfusion}. TinyFusion makes this dependence explicit by designing the pruned architecture for post-fine-tuning recoverability~\cite{fang2025tinyfusion}. In audio, Singh et al.~\cite{singh2026pruning} L1-prune the AudioLDM U-Net and then fine-tune the pruned model on AudioCaps for $10^{6}$ steps. Their recovered model substantially improves the reported in-domain metrics and can surpass the unpruned release in FAD.

These results establish that a pruned network can regain benchmark performance. Recent methods also make recovery progressively more central to the compression problem. Diff-Pruning uses diffusion-aware gradients to identify removable structure and then limits the retraining needed to obtain a useful model~\cite{fang2023diffpruning}. BK-SDM removes residual and attention blocks and transfers behavior into the smaller U-Net through feature distillation~\cite{kim2024bksdm}. TinyFusion goes further by selecting a shallow architecture for how well it is expected to perform \emph{after} fine-tuning, rather than only for its error immediately after pruning~\cite{fang2025tinyfusion}. Across these approaches, the final compressed model is therefore defined not only by what was removed, but also by what adaptation can recover.

That progress makes the meaning of ``recovery'' more important. Recovery fine-tuning is an adaptation procedure carried out on one data distribution and, for AudioLDM, one training duration. A recovered score can therefore mix two effects, restoration of damage introduced by pruning and specialization toward the recovery conditions. The matched dense model given the same recovery fine-tune is unavailable in the release studied here, so we do not attempt to separate those mechanisms. We instead ask the narrower evaluation question that can be answered from the released checkpoints.

\subsection{What should recovery evaluation measure?}
Text-audio similarity from CLAP is widely used as an automatic proxy for prompt alignment~\cite{wu2023clap,huang2023makeanaudio,ghosal2023tango}. FAD measures distributional similarity but is sensitive to sample size and embedding choice~\cite{kilgour2019fad,gui2024fad,chung2025kad}. Human-CLAP~\cite{takano2025humanclap}, event classifiers such as PANNs~\cite{kong2020panns} and frame-level grounding models such as FineLAP~\cite{li2026finelap} provide complementary views.

The evaluation problem is also changing as TTA systems mature. A single quality score is useful for ranking systems under a fixed recipe, but it says little about what happens when prompts, generation settings or use cases move away from that recipe. TTA-Bench makes this broader concern explicit by evaluating current TTA models along dimensions that include robustness and generalization~\cite{wang2026ttabench}. Our scope is narrower. We do not propose another general benchmark. We ask how the same principle should be applied when the object being evaluated is specifically the \emph{gain produced by recovery fine-tuning}.

No one metric, moreover, makes an absolute score intrinsically meaningful. A CLAP value can be high or low because of the model, the prompt battery or the scorer's own baseline. A useful recovery evaluation therefore separates three quantities. The final score describes the checkpoint that a user receives. The paired change from P to \PFT{} describes what fine-tuning added. The fraction of a reference gap closed by that change describes how large the addition is relative to a target such as the dense model or real audio. Keeping these quantities separate prevents a high final score from being interpreted automatically as evidence of broad recovery.

We evaluate the \emph{increment caused by recovery} and read it against two anchors. A shuffled-caption floor asks whether generated audio aligns with its own prompt better than with unrelated prompts. Real audio asks how much of the task-level gap remains even when the prompt is paired with the source recording. These anchors make the recovery gain easier to interpret than an isolated cosine value and distinguish a statistically resolved change from a practically substantial restoration.

\begin{figure*}[t]
\centering
\includegraphics[width=\textwidth]{figs/fig1_operating_points.pdf}
\caption{Recovery across the operating points measured in this study. Panels (a) and (b) show mean CLAP cosine for P and \PFT{} as requested duration changes, together with matched dense generations, real audio and each cell's chance floor. Whiskers show the $95\%$ CI of paired recovery gain $R$. Panel (c) compares in-domain AudioCaps with held-out hip-hop captions under matched duration. Panel (d) shows FineLAP frame-level grounding gain through the $10.24$\,s clip. Panel (e) compares separately generated short clips with short crops from native-duration generations.}
\label{fig:ops}
\end{figure*}

